\chapter{緒論}

\section{研究背景及意義}
緒論部分應說明選題背景、研究必要性、研究問題、研究動機與研究目標，並清楚區分已有研究工作與本人完成的工作。

本模板使用 \texttt{ref.bib} 管理參考文獻，例如可使用 \cite{lecun2015deep} 進行引用。

\subsection{選題背景}
每個圖均應包含圖序號與圖題，圖題置於圖下方；圖應在正文中先提及後出現。示例如 \cref{fig:format} 所示。

\begin{figure}[htbp]
  \centering
  \fbox{%
    \begin{minipage}[c][5.0cm][c]{0.75\textwidth}
      \centering
      圖片占位符\\[0.4cm]
      \small 請在此處替換為論文圖片。
    \end{minipage}
  }
  \caption{論文中圖的格式要求}
  \label{fig:format}
\end{figure}
